% Lembar 1 kuesioner penelitian Bab 16: butir X dan Y dengan skala frekuensi.
% Urutan butir di lembar: B1 = X1, B2 = Y1, B3 = X2, B4 = Y2, B5 = X3,
% B6 = Y3, B7 = X4, B8 = Y4, B9 = X5, B10 = Y5 (lihat Tabel 16.2 di buku).
\documentclass[11pt]{article}
% Gaya bersama untuk kedua lembar kuesioner penelitian Bab 16.
\usepackage[a4paper, top=1.6cm, bottom=1.6cm, left=2cm, right=2cm]{geometry}
\usepackage{fontspec}
\setmainfont[
  Path = ../fonts/,
  UprightFont = TitilliumWeb-Regular.ttf,
  BoldFont = TitilliumWeb-Bold.ttf,
  ItalicFont = TitilliumWeb-Italic.ttf,
  BoldItalicFont = TitilliumWeb-BoldItalic.ttf
]{TitilliumWeb}
\usepackage[bahasai]{babel}
\usepackage{xcolor}
\usepackage{graphicx}
\usepackage{tabularx}
\usepackage{array}
\usepackage{tikz}

\definecolor{praditagreen}{RGB}{37, 113, 71}
\setlength{\parindent}{0pt}
\setlength{\parskip}{0.5em}
\pagestyle{empty}

% Bagian yang harus diisi peneliti sebelum lembar dicetak.
\newcommand{\isian}[1]{\textcolor{red}{[#1]}}

% Kepala lembar: logo, judul penelitian, dan nomor lembar.
\newcommand{\kepala}[1]{%
  \begin{minipage}[c]{0.22\textwidth}
    \includegraphics[width=\linewidth]{../../figures/images/logo_universitas.png}
  \end{minipage}\hfill
  \begin{minipage}[c]{0.74\textwidth}
    \raggedleft
    {\Large\bfseries\color{praditagreen} Kuesioner Penelitian}\\[2pt]
    {\large Belajar Pemrograman dengan Bantuan AI}\\[2pt]
    {\bfseries #1}
  \end{minipage}
  \par\vspace{2pt}{\color{praditagreen}\rule{\textwidth}{1pt}}\par
}

\newcommand{\bagian}[1]{\par\medskip{\large\bfseries\color{praditagreen}#1}\par}

% Kotak dan lingkaran digambar dengan TikZ agar tidak bergantung pada font simbol.
\newcommand{\kotakpilih}{\tikz[baseline=-0.15em]\draw (0,0) rectangle (0.3,0.3);}
\newcommand{\lingkaran}{\tikz[baseline=-0.6ex]\draw (0,0) circle (0.17);}

% Pilihan jawaban: kotak tidak terpisah dari labelnya saat baris dipatah.
\newcommand{\opsi}[1]{\mbox{\kotakpilih~#1}\hspace{1em plus 0.5em}}
\newenvironment{barispilihan}{\par\raggedright}{\par}

% NIM ditulis di kedua lembar untuk memasangkannya. NIM diganti kode acak
% sebelum analisis dan tidak dimuat dalam publikasi.
\newcommand{\kotaknim}{%
  Tuliskan NIM Anda, satu angka di setiap kotak.

  \begin{tabular}{|*{10}{p{0.45cm}|}}
    \hline
    \rule{0pt}{0.8cm} & & & & & & & & & \\
    \hline
  \end{tabular}
}

% Tabel butir dengan lima pilihan lingkaran. Setiap baris diakhiri \pilihan.
\newenvironment{tabelskala}{%
  \renewcommand{\arraystretch}{1.45}%
  \tabularx{\textwidth}{>{\raggedright\arraybackslash}p{0.8cm}X*{5}{>{\centering\arraybackslash}p{0.65cm}}}
  \hline
  \textbf{No} & \textbf{Pernyataan} & \textbf{1} & \textbf{2} & \textbf{3} & \textbf{4} & \textbf{5} \\
  \hline
}{\endtabularx}
\newcommand{\pilihan}{& \lingkaran & \lingkaran & \lingkaran & \lingkaran & \lingkaran \\ \hline}

\newcommand{\garis}{\par\vspace{0.55cm}\noindent\rule{\textwidth}{0.4pt}}


\begin{document}

\kepala{Lembar 1 dari 2}

Kuesioner ini bagian dari penelitian tentang cara mahasiswa memakai AI saat belajar pemrograman dasar. Penelitian dilakukan oleh Alfa Yohannis dari Program Studi Informatika, Universitas Pradita.

\begin{itemize}
  \setlength{\itemsep}{0pt}
  \item Kuesioner terdiri atas dua lembar yang diisi pada hari berbeda. Lembar ini adalah Lembar 1.
  \item Tuliskan NIM Anda di kedua lembar. NIM dipakai untuk memasangkan kedua lembar dan menghubungkan jawaban Anda dengan nilai kuis minggu pertama.
  \item Lembar dikumpulkan oleh \isian{pengumpul data}, lalu NIM diganti kode acak sebelum data dianalisis. Hasil penelitian yang dipublikasikan tidak memuat NIM.
  \item Tidak ada jawaban benar atau salah. Jawablah sesuai dengan yang benar-benar Anda lakukan.
  \item Keikutsertaan bersifat sukarela dan tidak memengaruhi nilai mata kuliah. Anda boleh tidak mengisi atau berhenti kapan saja.
  \item Dengan mengisi dan menyerahkan lembar ini, Anda menyatakan bersedia ikut serta.
\end{itemize}

Pengisian membutuhkan sekitar 10 menit. Pertanyaan tentang penelitian ini dapat disampaikan kepada \isian{kontak peneliti}.

\bagian{NIM}
\kotaknim

\bagian{Bagian A. Pemakaian AI secara umum}

A1. Layanan AI apa yang paling sering Anda pakai untuk mata kuliah ini? Pilih satu.

\begin{barispilihan}
\opsi{ChatGPT}\opsi{Gemini}\opsi{Claude}\opsi{Microsoft Copilot}\opsi{DeepSeek}\opsi{Lainnya: \rule{2.5cm}{0.4pt}}\opsi{Tidak pernah memakai AI}
\end{barispilihan}

A2. Seberapa sering Anda memakai AI untuk mata kuliah ini selama semester ini? Pilih satu.

\begin{barispilihan}
\opsi{Tidak pernah}\opsi{Kurang dari sekali seminggu}\opsi{1 sampai 2 kali seminggu}\opsi{3 sampai 5 kali seminggu}\opsi{Hampir setiap hari}
\end{barispilihan}

\newpage
\bagian{Bagian B. Kebiasaan belajar dengan AI}

Seberapa sering Anda melakukan hal berikut saat belajar mata kuliah Pemrograman Dasar selama semester ini? Beri tanda silang pada satu lingkaran untuk setiap pernyataan.

Skala: 1 = Tidak pernah, 2 = Jarang, 3 = Kadang-kadang, 4 = Sering, 5 = Selalu.

\begin{tabelskala}
  B1 & Saya memakai \textit{prompt} terarah dari buku pegangan kuliah saat mempelajari listing atau mengerjakan latihan. \pilihan
  B2 & Saya meminta AI menuliskan kode lengkap untuk tugas atau latihan. \pilihan
  B3 & Saya meminta AI memberi pertanyaan untuk menguji pemahaman saya. \pilihan
  B4 & Saya menyalin kode dari AI tanpa memahami cara kerjanya. \pilihan
  B5 & Saya mencoba menjawab pertanyaan dari AI sebelum melihat penjelasannya. \pilihan
  B6 & Saat muncul pesan kesalahan, saya langsung bertanya ke AI tanpa mencoba memperbaikinya sendiri. \pilihan
  B7 & Saya meminta AI membuat soal latihan baru, lalu mengerjakannya sendiri. \pilihan
  B8 & Saya merasa kesulitan ketika harus menulis program tanpa bantuan AI. \pilihan
  B9 & Saya menjalankan sendiri kode dari penjelasan AI untuk memeriksa kebenarannya. \pilihan
  B10 & Tugas yang saya kumpulkan sebagian besar kodenya berasal dari AI. \pilihan
\end{tabelskala}

\bagian{Bagian C. Pertanyaan terbuka}

Ceritakan secara singkat cara Anda biasanya memakai AI saat belajar mata kuliah ini. Bagian ini boleh dikosongkan.
\garis\garis\garis

\medskip
Terima kasih. Lembar 2 akan dibagikan pada \isian{tanggal pengisian Lembar 2}.

\end{document}
